\section{Methodology}

The framework maps temporally indexed enrollment data to weekly hazard estimates, simulated regime contrasts, and subgroup diagnostics. It combines person-period construction, leakage-aware hazard modeling, policy simulation, and subgroup analysis under explicit inferential limits.

\subsection{Motivation, Contributions, and Research Questions}

\paragraph*{Motivation.}
Static risk models primarily answer \textit{who} is at risk, but they do not by themselves determine \textit{when to act} under an auditable operational rule \cite{Prenkaj2021SurveyMLDropout,AndradeGiron2023DropoutMLReview}. In observational settings, the comparison of intervention scenarios must therefore be framed as a structural, model-implied exercise rather than as a causally identified effect estimate \cite{Connolly2018TrialsEvidenceBased,Gallagher2020ChallengeBased,Sidiropoulos2022AgencySustainability,Igelstrom2022CausalInferenceObs,DoutreligneVaroquaux2025PredictiveDecisionCausal}. This study is motivated by that gap: the need for a temporally explicit framework that links weekly risk estimation to executable policy simulation and subgroup-sensitive interpretation.

\paragraph*{Contributions.}
The framework has four components. First, it formalizes a discrete-time hazard pipeline for estimating survival trajectories from temporally indexed student records. Second, it introduces a policy simulation layer that compares \textit{shock} and \textit{stateful mechanism-aware} regimes under a common intervention rule. Third, it adds a subgroup fairness layer to quantify how the same policy changes outcome gaps across observable groups. Fourth, it specifies an executable protocol grounded in the notebook implementation and exported artifacts.

\paragraph*{Research Questions.}
RQ1 asks whether the temporal model produces weekly hazard estimates that are sufficiently discriminative and calibrated to support decision-making. RQ2 asks what structural survival contrast emerges between simulated regimes under an explicit policy rule. RQ3 asks whether the same policy changes subgroup gaps, and with what uncertainty.

\paragraph*{Propositions.}
\begin{proposition}[Temporal bridge]
If $\hat h_{it}$ is estimated on person-period data without temporal leakage, then $\hat S_i(t)$ defines candidate temporal risk windows for downstream analysis.
\end{proposition}

\begin{proposition}[Regime contrast]
Under a fixed policy contract,
\begin{equation*}
\Delta S(t)=\bar S^{(1)}(t)-\bar S^{(0)}(t)
\end{equation*}
defines a structural contrast between simulated regimes.
\end{proposition}

\begin{proposition}[Subgroup heterogeneity]
When the same policy is applied to observable groups,
\begin{equation*}
\Delta Gap(t)=Gap^{(1)}(t)-Gap^{(0)}(t)
\end{equation*}
quantifies the differential change in trajectories under the same simulation engine.
\end{proposition}

\subsection{Notation and Evaluation Horizons}

Table~\ref{tab:method_notation} summarizes the notation used throughout the methodological tracks.

\begin{table}[!t]
\centering
\footnotesize
\renewcommand{\arraystretch}{1.05}
\begin{tabular}{@{}p{0.25\columnwidth} p{0.67\columnwidth}@{}}
\hline
Symbol & Definition \\
\hline
$i$ & enrollment (unit of analysis) \\
$t$ & discrete week \\
$X_{it}$ & observed covariates at time $t$ \\
$E_i$ & primary event indicator \\
$C_i$ & censoring time \\
$\hat h_{it}$ & predicted weekly hazard \\
$\hat S_i(t)$ & predicted survival \\
$\bar S^{(a)}(t)$ & mean survival under regime $a$ \\
$\Delta S(t)$ & structural survival contrast \\
$Gap^{(a)}(t)$ & between-group gap under regime $a$ \\
$\Delta Gap(t)$ & change in gap between regimes \\
$T_{policy}$ & primary substantive reporting horizon \\
$T_{eval\_policy}$ & temporal support evaluated in policy simulation \\
$T_{eval\_metrics}$ & stable horizon for IPCW-based metrics \\
$g_{min}$ & minimum stability threshold for $\hat G(t)$ \\
\hline
\end{tabular}
\caption{Main notation used throughout the Methodology section.}
\label{tab:method_notation}
\end{table}

The protocol distinguishes between substantive and technical horizons:
\begin{equation*}
T_{eval\_metrics}=\max\{t:\hat G(t)\ge g_{min}\},
\end{equation*}
where $T_{policy}$ denotes the main reporting horizon and $T_{eval\_policy}$ denotes the temporal support over which simulated regime trajectories are evaluated. Policy interpretation and censoring-weighted metric stability do not necessarily operate on the same horizon.

In the current exported protocol, these horizons are fixed as $T_{policy}=18$, $T_{eval\_policy}=38$, and $T_{eval\_metrics}=37$, with $g_{min}=0.05$ recorded as the operational stability threshold in the evidence package.

\subsection{Analytical Tracks}

The framework spans the bridge
\texttt{temporal hazard $\rightarrow$ survival trajectories $\rightarrow$}
\texttt{policy simulation $\rightarrow$ subgroup diagnostics}.
The implemented tracks are:
\begin{enumerate}
\item temporal pipeline and data construction;
\item leakage prevention and temporal splitting;
\item censoring and IPCW;
\item primary discrete-time hazard model (RQ1);
\item benchmark model family;
\item policy simulation (RQ2);
\item subgroup fairness diagnostics (RQ3).
\end{enumerate}

The framework does not identify causal effects from observational exposure. It defines the modeling and simulation machinery under which later contrasts are generated and interpreted.

\subsection{Problem Setup and Unit of Analysis}

The unit of analysis is enrollment $i$ observed over discrete weekly time $t$:
\begin{equation*}
i \in \{1,\dots,N\}, \quad t \in \{0,\dots,T\}, \quad X_{it}\in\mathbb{R}^p.
\end{equation*}

The analytical table is constructed in person-period format, with one row per enrollment-week. This representation is required because the target quantity is not a single static outcome per student, but a time-indexed conditional event risk that evolves over the observed trajectory.

\subsection{Primary Endpoint and Censoring}

The primary event is administrative withdrawal, operationalized as \texttt{Withdrawn} with a valid \texttt{date\_unregistration}. Cases labeled \texttt{Withdrawn} without a valid date are treated as censored under the primary endpoint definition.

\begin{equation*}
\begin{aligned}
E_i=\mathbb{1}\{&\texttt{final\_result}_i=\texttt{Withdrawn} \\
&\land\ \texttt{date\_unregistration}_i\ \text{is valid}\}.
\end{aligned}
\end{equation*}

\begin{equation*}
t_i^{event}=\left\lfloor \frac{\texttt{date\_unregistration}_i}{7} \right\rfloor \quad (E_i=1)
\end{equation*}

\begin{equation*}
t_i^{final}=\begin{cases}
t_i^{event}, & E_i=1, \\
t_i^{last\_obs}, & E_i=0.
\end{cases}
\end{equation*}

For non-events, $t_i^{last\_obs}$ is operationalized as the last observed VLE week and should be interpreted as an observation proxy rather than an administrative censoring timestamp.

Events are counted only when administrative timing is observed and compatible with the weekly survival setup.

\subsection{Person-Period Construction and Terminal Label}

Each enrollment is expanded weekly up to $t_i^{final}$, yielding one terminal event label at most:
\begin{equation*}
event_{it}=\mathbb{1}\{E_i=1 \land t=t_i^{final}\}
\end{equation*}

\begin{equation*}
\sum_{t=0}^{t_i^{final}} event_{it} \le 1.
\end{equation*}

This construction yields a longitudinal risk set for each enrollment and, if applicable, a single terminal event week.

\subsection{Temporal Covariate Engineering}

Dynamic covariates are defined on a weekly grid under temporally safe logic:
\begin{equation*}
w(d)=\max\left(\left\lfloor d/7 \right\rfloor,0\right)
\end{equation*}

\begin{equation*}
A_{it}=\mathbb{1}\{\texttt{total\_clicks}_{it}>0\}
\end{equation*}

\begin{equation*}
Recency_{it}=\begin{cases}
0, & A_{it}=1, \\
Recency_{i,t-1}+1, & A_{it}=0
\end{cases}
\end{equation*}

\begin{equation*}
Streak_{it}=\begin{cases}
Streak_{i,t-1}+1, & A_{it}=1, \\
0, & A_{it}=0.
\end{cases}
\end{equation*}

The main feature set includes \texttt{total\_clicks}, \texttt{recency}, \texttt{streak}, \texttt{submitted\_this\_week}, and contextual covariates. These features encode short-term engagement dynamics while preserving temporal ordering.

\subsection{Stratified Temporal Split and Leakage Prevention}

The train-test split is performed at the enrollment level, stratified jointly by event status and temporal bucket:
\begin{equation*}
time\_for\_split_i=\begin{cases}
t_i^{event}, & E_i=1, \\
t_i^{final}, & E_i=0
\end{cases}
\end{equation*}

\begin{equation*}
stratum_i=(E_i,\ bucket_q(time\_for\_split_i)).
\end{equation*}

The implemented configuration uses $q=4$ and \texttt{test\_size}=0.30. Calibration and validation procedures further use enrollment-grouped folds via GroupKFold.

The design prevents leakage across rows belonging to the same enrollment and preserves temporal composition across train and test partitions. The grouped validation logic follows the broader concern with leakage in structured machine-learning evaluation and cross-validation for dependent data \cite{KapoorNarayanan2023LeakagePatterns,RobertsEtAl2017CrossValidationEcography}.

\subsection{Primary Discrete-Time Hazard Model (RQ1)}

Following standard discrete-time event-history formulations \cite{SingerWillett1993DiscreteTimeSurvival,Allison1982DiscreteTimeMethods}, the primary modeling track estimates weekly conditional hazard and derives survival by cumulative product:
\begin{equation}
\label{eq:hazard_weekly}
h_{it}=P(event_{it}=1 \mid T_i\ge t, X_{it})
\end{equation}

\begin{equation*}
\text{logit}(\hat h_{it})=\beta_0+\beta^\top X_{it}
\end{equation*}

In implementation, the raw weekly probabilities are post-hoc calibrated through grouped sigmoid calibration (Platt scaling) to support interpretable hazard probabilities under the same enrollment-respecting split logic \cite{Platt1999ProbabilisticOutputs}.

\begin{equation}
\label{eq:survival_product}
\hat S_i(t)=\prod_{k\le t}(1-\hat h_{ik})
\end{equation}

The primary track yields weekly risk trajectories rather than a single end-of-course classification score. This layer underlies RQ1 and the downstream policy simulation.

\subsection{Censoring Track and IPCW}

The censoring component combines estimation of censoring survival with truncated inverse probability of censoring weights:
\begin{equation*}
\hat G_{it}=P(C_i\ge t \mid X_{it})
\end{equation*}

\begin{equation*}
w_{it}=\frac{1}{\max(\hat G_{it}, g_{min})}.
\end{equation*}

This weighting logic follows standard survival-prediction evaluation under right censoring \cite{Graf1999AssessmentSurvivalPrediction,Gerds2006ConsistentBrier}. The protocol separates the policy reporting horizon from the technical horizon used for censoring-weighted metrics because IPCW evaluation becomes unstable when $\hat G(t)$ is too small, even if simulated regime trajectories remain computable beyond that point. In the current protocol, $g_{min}=0.05$ yields $T_{eval\_metrics}=37$, while $T_{eval\_policy}=38$ is retained only as the raw support horizon for trajectory reporting.

Because non-event follow-up is anchored to last observed VLE activity, IPCW validity in this setup relies on conditional independent censoring given $X_{it}$ rather than unconditional random censoring. A sensitivity probe exported in \texttt{table\_censoring\_anchor\_sensitivity.csv} shows that across anchor variants the censoring-model test AUC ranges from $0.9155$ to $0.9299$, while capped-weight share in test ranges from $0.4945$ to $0.6576$ (current anchor at the upper end).

\subsection{Policy Simulation (RQ2): Shock vs. Mechanism-Aware}

The current policy contract separates external anchoring from internal scenario assumptions. The operational intervention rule is anchored to Kay and Bostock \cite{KayBostock2023PowerNudge}: ``flag if no LMS engagement in last 7 days ($\approx \text{recency} \ge 1$ week)'', with weekly checking frequency and trigger parameter $r^*=1$. This operational anchor is recorded in \texttt{table\_policy\_spec.csv}.

The active window is fixed at $W_{\text{weeks}}=2$, consistent with a short-term digital nudge framing \cite{KayBostock2023PowerNudge}. In contrast, effect intensity is treated as a scenario parameter because the observational setting does not identify a unique causal magnitude. The mechanism-aware shared schedule is frozen as $\alpha_{\text{week0}}=0.35$, $\alpha_{\text{week1}}=0.10$, \nolinkurl{decay_type=kb2023_step_2w}, and \nolinkurl{window_exclusive_upper=True}. For the shock regime, \nolinkurl{delta_shock_ablation} is reported as three scenario intensities: $0.08$ (anchored conservative scenario, qualitatively aligned with moderate intervention effects synthesized by Sneyers and De Witte \cite{SneyersDeWitte2018InterventionsMeta}), $0.20$ (hypothetical A), and $0.60$ (hypothetical B stress test).

The intervention rule is triggered by recency crossing a threshold within an activation window:
\begin{equation*}
t_i^*=\min\{t: Recency_{it}\ge r^*\}
\end{equation*}

\begin{equation*}
active_{it}=\mathbb{1}\{t\in[t_i^*,t_i^*+W)\}.
\end{equation*}

Under the \textit{shock} regime, the policy acts directly on the baseline hazard:
\begin{equation*}
\hat h^{(1,shock)}_{it}=\begin{cases}
\hat h^{(0)}_{it}(1-\delta_{shock}), & active_{it}=1, \\
\hat h^{(0)}_{it}, & active_{it}=0.
\end{cases}
\end{equation*}

Under the \textit{mechanism-aware} regime, the policy propagates through counterfactual covariate updates:
\begin{equation*}
\hat h^{(1,mech)}_{it}=f(X^{cf}_{it}).
\end{equation*}

In the mechanism-aware regime, $X^{cf}_{it}$ is updated statefully across weeks. In the exported run, the strongest perturbation channel is click-based engagement (\texttt{total\_clicks}), used as an LMS proxy analogous to participation and re-engagement dynamics discussed by Kay and Bostock \cite{KayBostock2023PowerNudge}, rather than as a literal replication of the source study variables. The shock regime applies an immediate model-level perturbation, whereas the mechanism-aware regime changes the evolving covariate path through which future hazards are generated.

Operational details of this mechanism-aware operator are exported explicitly in \nolinkurl{table_policy_mech_operator_spec.csv} (channel mode, timing, overwrite rule, and propagation rule). The schedule sensitivity battery is exported in \nolinkurl{table_rq2_sensitivity_grid.csv} and varies trigger/schedule parameters while preserving the same plug-in modeling stack used for baseline scoring.

The updated framework also keeps mechanism-aware execution auditable through run-level diagnostics exported as \nolinkurl{table_policy_activation_summary.csv}, \nolinkurl{table_policy_covariate_overwrites.csv}, and \nolinkurl{table_policy_covariate_propagation_checks.csv}, alongside the operator contract.

\subsection{Structural Survival Contrast (RQ2)}

The central policy estimand is mean survival under each regime:
\begin{equation*}
\bar S^{(a)}(t)=\frac{1}{N}\sum_i \hat S_i^{(a)}(t)
\end{equation*}

\begin{equation}
\label{eq:delta_s}
\Delta S(t)=\bar S^{(1)}(t)-\bar S^{(0)}(t).
\end{equation}
Positive values indicate a survival gain under the policy regime relative to baseline.

The interpretation of $\Delta S(t)$ is that of a \textit{simulated regime contrast} implied by the fitted model and the policy contract. It is not interpreted as a causally identified treatment effect \cite{Igelstrom2022CausalInferenceObs,DoutreligneVaroquaux2025PredictiveDecisionCausal}.

\subsection{Fairness and Subgroup Track (RQ3)}

The same policy is evaluated across observable subgroups through a gap-based contrast:
\begin{equation*}
\bar\mu_g^{(a)}(t)=\frac{1}{N_g}\sum_{i:g(i)=g}\hat S_i^{(a)}(t)
\end{equation*}

\begin{equation*}
Gap^{(a)}(t)=\bar\mu_1^{(a)}(t)-\bar\mu_0^{(a)}(t)
\end{equation*}

\begin{equation}
\label{eq:delta_gap}
\Delta Gap(t)=Gap^{(1)}(t)-Gap^{(0)}(t).
\end{equation}
The sign of $\Delta Gap(t)$ indicates whether the policy enlarges or reduces the between-group difference at the horizon of interest. Uncertainty is estimated via bootstrap, consistent with subgroup-oriented fairness auditing that reports performance and calibration gaps with interval estimates \cite{PessachShmueli2022FairnessReview,Lu2022ReliabilityFairnessAudit,Opoku2025AccuracyFairness}.

\subsection{Benchmark Track}

The benchmark track provides model-family reference comparisons without replacing the primary hazard pipeline. In the current main-text reporting, the benchmark centers on the temporal hazard model, its IPCW-weighted variant, and a non-temporal Random Survival Forest baseline at aligned horizons \cite{MogensenIshwaranGerds2012RSF}. Additional model families, including Cox time-varying specifications \cite{Cox1972RegressionModels}, can serve as auxiliary references.

\subsection{Operational Assumptions}

Interpretation rests on the following assumptions:
\begin{itemize}
\item the event and censoring definitions correctly represent the observational endpoint;
\item covariates used at week $t$ are available up to week $t$ only, with no temporal leakage;
\item policy simulation defines a model-implied structural contrast rather than a causally identified intervention effect;
\item the stability of censoring-weighted metrics is evaluated at the technical horizon $T_{eval\_metrics}$.
\end{itemize}

\subsection{Interpretive Scope}

The reported results for $\Delta S$ and $\Delta Gap$ are \textit{model-implied simulated regime contrasts}, not causally identified estimates of intervention effects.

\subsection{Reproducibility}

Code for preprocessing, temporal splitting, training and calibration, policy simulation, subgroup fairness analysis, and artifact export (\texttt{outputs\_v2}) is available at
\url{https://github.com/rafa-rodriguess/TCM-Student-Dropout}
\cite{RafaRepoTCMStudentDropout}. The reference dataset used in the pipeline is OULAD \cite{Kuzilek2017OULAD}.

The artifact package records the full policy
contract in \nolinkurl{table_policy_spec.csv}, the scenario catalog in
\nolinkurl{table_policy_scenarios_main.csv} and
\nolinkurl{table_policy_scenario_params.csv}, and scenario-resolved trajectory
outputs in \nolinkurl{table_policy_deltaS_by_week_by_scenario.csv} and
\nolinkurl{table_policy_deltaS_at_horizons_by_scenario.csv}. It also records
the dual-horizon semantics in \nolinkurl{table_policy_horizons_dual.csv}, the mechanism-aware operator contract in \nolinkurl{table_policy_mech_operator_spec.csv}, and the RQ2 schedule sensitivity grid in \nolinkurl{table_rq2_sensitivity_grid.csv}. Under the current exported protocol, these artifacts jointly encode the
anchored trigger reference, the frozen shared mechanism-aware schedule, the
anchored-vs-hypothetical shock intensity split, and the distinction between
$T_{\text{policy}}=18$, $T_{\text{eval,policy}}=38$, and
$T_{\text{eval,metrics}}=37$.